\documentclass[10pt]{article}
\usepackage{float}
\usepackage{tabularx}
\usepackage{mathtools}
\usepackage{amsmath, amssymb, amsthm}
\usepackage{enumitem}
\usepackage{geometry}
\usepackage{fancyhdr}
\usepackage{titlesec}
\usepackage{tikz}
\usetikzlibrary{arrows.meta, positioning, calc, shapes.geometric, fit}
\usepackage{hyperref}
\usepackage{bookmark}
\geometry{margin=1in}

\pagestyle{fancy}
\fancyhf{}
\title{SFWRENG 3MX3 2025 Midterm 2 Solutions}
\author{Matthew Farah, \href{mailto:farahm11@mcmaster.ca}{$\langle$farahm11@mcmaster.ca$\rangle$}, 400468251}
\date{\today}

\hypersetup{
	colorlinks=true,
	linkcolor=blue,
	filecolor=magenta,
	urlcolor=blue,
	citecolor=blue,
	anchorcolor=blue
}

\tikzset{
	block/.style = {draw, rectangle, minimum height=1cm, minimum width=1.4cm},
	sum/.style   = {draw, circle, inner sep=5pt},
	gain/.style  = {draw, isosceles triangle, isosceles triangle apex angle=60,
		minimum height=1cm, minimum width=1cm, shape border rotate=180},
	connector/.style={thick},
}

\DeclareMathOperator{\sinc}{sinc}

\fancyhead{}
\fancyhead[L]{Matthew Farah, \href{mailto:farahm11@mcmaster.ca}{$\langle$farahm11@mcmaster.ca$\rangle$}, 400468251}
\fancyhead[R]{SFWRENG 3MX3 2025 Midterm 2 Solutions}

\setlength{\parindent}{0cm}

\newcounter{question}
\renewcommand{\thequestion}{\arabic{question}}

\newenvironment{question}[1][]{
	\newpage
	\refstepcounter{question}
	\phantomsection
	\addcontentsline{toc}{section}{\protect{\large{Question \thequestion} - #1}}
	\vspace{1em}
	\par
	\large\textbf{Question \thequestion}
	\vspace{.2cm}
	\par
	\setcounter{subquestion}{0}
	\setcounter{step}{0}
}{\vspace{.5em}}

\newenvironment{solution}{
	\vspace{1em}\par\large\textbf{{Solution:}}\par\vspace{1em}
}{
	\vspace{1em}
}

\newcounter{subquestion}
\renewcommand{\thesubquestion}{\alph{subquestion}}

\newenvironment{subquestion}{
	\refstepcounter{subquestion}
	\vspace{1em}\par\large\textbf{(\thesubquestion)}\hspace{-.5em}
	\setcounter{step}{0}
	\phantomsection
	\addcontentsline{toc}{subsection}{\protect{\large{\thequestion.\thesubquestion}}}
}{
	\par
	\vspace{1em}
}

\makeatother
\makeatletter

\newcounter{step}
\renewcommand{\thestep}{\Roman{step}}

\newenvironment{step}{
	\refstepcounter{step}
	\vspace{1.5em}\par\noindent\large\textbf{(\thestep)}
}{
	\par
	\vspace{1.5em}
}

\makeatother

\newcolumntype{Y}{>{\centering\arraybackslash}X}
\renewcommand{\arraystretch}{1.8}

\fontsize{10pt}{14pt}\selectfont

\begin{document}

\maketitle

\tableofcontents
\newpage

\begin{question}[Block Diagrams]
	Given the following block diagram of a discrete system, find its frequency response.

	\begin{center}
		\begin{tikzpicture}
			\node[sum] (sum1) {$+$};
			\node[sum, right = 2cm of sum1] (sum2) {$+$};
			\node[sum, right = 4cm of sum2] (sum3) {$+$};

			\node[gain, above left = 1.5cm and .5cm of sum1, shape border rotate = 0] (g1) {$.1$};
			\node[gain, right = 1.5cm of sum2, shape border rotate = 0] (g2) {$.3$};
			\node[gain, below right = 1.5cm and .5cm of sum2] (g3) {$.5$};
			\node[gain, below right = 4cm and 1cm of sum2, shape border rotate = 0] (g4) {$.2$};

			\node[block, below right = 1.3cm and 2cm of sum2] (d1) {$D$};
			\node[block, below = 3.625cm of sum2] (d2) {$D$};

			\draw[connector] (sum1) -- (sum2);
			\draw[connector] (sum2) -- (g2);
			\draw[connector] (g2) -- (sum3);

			\draw[connector] ($(sum3) - (1cm, 0)$) |- (d1);
			\draw[connector] (d1) -- (g3);
			\draw[connector] (g3) -| (sum2);

			\draw[connector] ($(sum1) - (2cm, 0)$) |- (d2);
			\draw[connector] (d2) -- (g4);
			\draw[connector] (g4) -| (sum3);

			\draw[connector] ($(sum1) - (2cm, 0)$) |- (g1);
			\draw[connector] (g1) -| (sum1);

			\draw[connector] (sum1) -- ++(-3cm,0);
			\draw[connector] (sum3) -- ++(3cm,0);
		\end{tikzpicture}
	\end{center}
\end{question}

\begin{solution}
	First, we simplify the block diagram by computing the frequency response, $H(\omega)$, of the indicated portion of the block diagram:

	\begin{center}
		\begin{tikzpicture}
			\node[sum] (sum1) {$+$};
			\node[sum, right = 2cm of sum1] (sum2) {$+$};
			\node[sum, right = 4cm of sum2] (sum3) {$+$};

			\node[gain, above left = 1.5cm and .5cm of sum1, shape border rotate = 0] (g1) {$.1$};
			\node[gain, right = 1.5cm of sum2, shape border rotate = 0] (g2) {$.3$};
			\node[gain, below right = 1.5cm and .5cm of sum2] (g3) {$.5$};
			\node[gain, below right = 4cm and 1cm of sum2, shape border rotate = 0] (g4) {$.2$};

			\node[block, below right = 1.3cm and 2cm of sum2] (d1) {$D$};
			\node[block, below = 3.625cm of sum2] (d2) {$D$};

			\draw[connector] (sum1) -- (sum2);
			\draw[connector] (sum2) -- (g2);
			\draw[connector] (g2) -- (sum3);

			\draw[connector] ($(sum3) - (1cm, 0)$) |- (d1);
			\draw[connector] (d1) -- (g3);
			\draw[connector] (g3) -| (sum2);

			\draw[connector] ($(sum1) - (2cm, 0)$) |- (d2);
			\draw[connector] (d2) -- (g4);
			\draw[connector] (g4) -| (sum3);

			\draw[connector] ($(sum1) - (2cm, 0)$) |- (g1);
			\draw[connector] (g1) -| (sum1);

			\draw[connector] (sum1) -- ++(-3cm,0);
			\draw[connector] (sum3) -- ++(3cm,0);

			\node[draw, inner xsep = .5cm, fit = (sum2) (d1) (g3) (g2), label = $H_1(\omega)$] {};
		\end{tikzpicture}
	\end{center}

	Since this is a simple feedback loop, we can determine that:

	\begin{align*}
		H_1(\omega) &= \frac{\frac{3}{10}}{1 - \left( \frac{3}{10} \right)\left(\frac{1}{2} \right)e^{-i\omega}} \\
		&= \frac{3}{10 - \frac{3}{2}e^{-i\omega}} \\
	\end{align*}

	Note here that the delay of one sample corresponds to $e^{-i\omega}$.

	We can now see that the block diagram has been simplified down to the following, once again, I have indicated the next part of the diagram to simplify:

	\begin{center}
		\begin{tikzpicture}
			\node[sum] (sum1) {$+$};
			\node[sum, right = 6cm of sum1] (sum3) {$+$};

			\node[gain, above left = 1.5cm and .5cm of sum1, shape border rotate = 0] (g1) {$.1$};
			\node[gain, below right = 4cm and 3cm of sum1, shape border rotate = 0] (g4) {$.2$};

			\node[block, right = 3cm of sum1] (H1) {$H_1(\omega)$};
			\node[block, below right = 3.79cm and 1cm of sum1] (d2) {$D$};

			\draw[connector] (sum1) -- (H1);
			\draw[connector] (H1) -- (sum3);

			\draw[connector] ($(sum1) - (2cm, 0)$) |- (d2);
			\draw[connector] (d2) -- (g4);
			\draw[connector] (g4) -| (sum3);

			\draw[connector] ($(sum1) - (2cm, 0)$) |- (g1);
			\draw[connector] (g1) -| (sum1);

			\draw[connector] (sum1) -- ++(-3cm,0);
			\draw[connector] (sum3) -- ++(3cm,0);

			\node[draw, inner xsep = 1cm, inner ysep = .5cm, fit = (sum1) (g1), label = $H_2(\omega)$] {};
		\end{tikzpicture}
	\end{center}

	\begin{align*}
		H_2(\omega) &= \frac{1}{10} + 1 \\
		&= \frac{11}{10} \\
	\end{align*}

	And we continue on in the same manner as before, giving us the frequency response of the entire system.

	\begin{center}
		\begin{tikzpicture}
			\node[sum, right = 6cm of sum1] (sum3) {$+$};

			\node[block, right = 3cm of sum1] (H1) {$H_1(\omega)$};
			\node[block] (H2) {$H_2(\omega)$};

			\node[gain, below right = 4cm and 3cm of H2, shape border rotate = 0] (g4) {$.2$};

			\node[block, below right = 3.79cm and 1cm of H2] (d2) {$D$};

			\draw[connector] (H1) -- (H2);
			\draw[connector] (H1) -- (sum3);

			\draw[connector] ($(sum1) - (2cm, 0)$) |- (d2);
			\draw[connector] (d2) -- (g4);
			\draw[connector] (g4) -| (sum3);

			\draw[connector] (H2) -- ++(-3cm,0);
			\draw[connector] (sum3) -- ++(3cm,0);

			\node[draw, inner xsep = 2cm, inner ysep = .5cm, fit = (sum3) (g4) (H1) (H2) (d2), label = $H(\omega)$] {};
		\end{tikzpicture}
	\end{center}

	\begin{align*}
		H(\omega) &= H_1(\omega)H_2(\omega) + \frac{1}{5}e^{-i\omega} \\
		&= \frac{3}{10 - \frac{3}{2}e^{-i\omega}} \cdot \frac{11}{10} + \frac{1}{5}e^{-i\omega} \\
		&= \frac{33}{100 - 15e^{-i\omega}} + \frac{1}{5}e^{-i\omega}
	\end{align*}
\end{solution}

\begin{question}[The Fourier Transform]
	Given is a filter by its impulse response: 

	\begin{align*}
		h(t) &=
		\begin{cases}
			1, \; \text{if} \; -1 \le t \le 0 \\
			-1, \; \text{if} \; 1 \le t \le 2 \\
			0, \; \text{otherwise}
		\end{cases} \\
	\end{align*}

	Compute the frequency response, $H(\omega)$ of the system.
\end{question}

\begin{solution}
	We know that the frequency response

	\begin{align*}
		H_0(\omega) &=
		\begin{cases}
			\frac{1}{2a}, \; \text{if} \; -1 \le t \le a \\
			0, \; \text{otherwise}
		\end{cases} \\
	\end{align*}

	corresponds to a signal whose impulse response is given by $\sinc(a\omega)$. Thus, by considering manipulations of $H_0(\omega)$, we can find the frequency response of the system without having to directly compute it.

	\begin{step}
		Find the frequency response for $-1 \le t \le 0$.
	\end{step}

	If we set $a = \frac{1}{2}$, and advance the impulse response by half a second, we would get an impulse response corresponding to $h(t) = \begin{cases} 1, \; \text{if} \; -1 \le t \le 0 \\ 0, \; \text{otherwise} \end{cases}$, thus giving that:

	\begin{align*}
		H_2(\omega) &= \sinc \left( \frac{1}{2}\omega \right)e^{\frac{1}{2}i\omega} \\
	\end{align*}

	\begin{step}
		Find the frequency response for $-1 \le t \le 2$.
	\end{step}

	Similarly, if we select $a = \frac{1}{2}$, delay the impulse response by 3 and a half seconds, and then scale by a factor of $-1$, we get the impulse response corresponding to $h(t) = \begin{cases} -1, \; \text{if} \; 1 \le t \le 2 \\ 0, \; \text{otherwise} \end{cases}$, thus giving that:

	\begin{align*}
		H_3(\omega) &= -\sinc \left( \frac{1}{2}\omega \right)e^{-\frac{3}{2}i\omega} \\
	\end{align*}

	\begin{step}
		Find the frequency response of the system.
	\end{step}

	Thus, putting $H_1(\omega)$ and $H_2(\omega)$ together, we find that:

	\begin{align*}
		H(\omega) &= H_1(\omega) + H_2(\omega) \\
		&= \sinc \left( \frac{1}{2}\omega \right)e^{\frac{1}{2}i\omega} - \sinc \left( \frac{1}{2}\omega \right)e^{-\frac{3}{2}i\omega}
	\end{align*}
\end{solution}

\begin{question}[Short Answer Questions]
	Answer the following questions:
\end{question}

\begin{subquestion}
	Compute the period of the signal given by $x(n) = \cos \left( \frac{5\pi}{4}n \right)$
\end{subquestion}

\begin{solution}
	The period, $P$ of $x(n)$ is the smallest value such that $\frac{5\pi}{4}P$ is a multiple of $2\pi$ where $P \in \mathbb{N}$. Thus, $P = 8$.
\end{solution}

\begin{subquestion}
	Compute the Fourier series coefficients ($X_k$) of the signal given by $x(t) = 3 + 2\cos(t) + \sin(3t)$.
\end{subquestion}

\begin{solution}
	By correlating the complex exponential expansion of each term in the signal, we arrive at the following Fourier series coefficients:

	\begin{align*}
		X_{-3} = \frac{-1}{2i} \quad & \quad X_{3} = \frac{1}{2i} \\
		X_{-2} = 0 \quad & \quad X_{2} = 0 \\
		X_{-1} = 1 \quad & \quad X_{1} = 1 \\
		X_0 &= 3 \\
	\end{align*}

	Note here that no signal has a frequency of $3$, so $X_{-2} = X_{2} = 0$.
\end{solution}

\begin{subquestion}
	How fast must one sample to observe something occurring at a frequency of 10Hz?
\end{subquestion}

\begin{solution}
	The sampling frequency must be \underline{strictly} greater than twice the frequency of what we wish to observe. Thus, we must sample at a rate of $>20$Hz.
\end{solution}

\begin{subquestion}
	Compute the discrete Fourier transform of the signal $x(n) = \delta(n)$.
\end{subquestion}

\begin{solution}
	The formula for the discrete Fourier transform is given by:

	\begin{align*}
		X_k &= \sum_{n = 0}^{N - 1}x(n)e^{-i\omega_0{n}{k}} \\
	\end{align*}

	Therefore, we find that:

	\begin{align*}
		X_k &= \sum_{n = 0}^{N - 1}x(n)e^{-i\omega_0{n}{k}} \\
		&= \sum_{n = 0}^{N - 1}\delta(n)e^{-i\omega_0{n}{k}} \\
		&= \delta(0)e^{-i\omega_0{0}{k}} \\
		&= e^{0} \\
		&= 1
	\end{align*}
\end{solution}

\begin{question}[Modulation]
	Given the signal $y(t) = x(t)\cos \left( \omega_c{t} \right)$. Given that the Fourier transform of $x(t)$ is $X(\omega)$, compute $Y(\omega)$.
\end{question}

\begin{solution}
	Note that multiplication in the time domain corresponds to convolution in the frequency domain. Thus, to find $y(t)$, we must find the frequency response of $\cos(\omega_c{t})$ and convolute that with $X(\omega)$. Let $H(\omega)$ be the Fourier transform of $\cos(\omega_{c}{t})$

	\begin{step}
		Compute $H(\omega)$.
	\end{step}

	To find $H(\omega)$, we must compute the continuous Fourier transform of $\cos \left( \omega_{c}{t} \right)$.

	\begin{align*}
		H(\omega) &= \int_{-\infty}^{\infty}\cos \left( \omega_{c}{t} \right)e^{-i\omega{t}}dt \\
		&= \frac{1}{2}\int_{-\infty}^{\infty} \left( e^{i\omega_{c}t} + e^{-i\omega_{c}t} \right)e^{-i\omega{t}}dt \\
		&= \frac{1}{2}\int_{-\infty}^{\infty} e^{it\omega(\omega_c - \omega)} + e^{-it(-\omega_c - \omega)}dt \\
	\end{align*}

	Now, using the fact that $\int_{-\infty}{\infty}e^{i\alpha{t}}dt = 2\pi\delta(a)$, we can conclude that:

	\begin{align*}
		H(\omega) &= \frac{1}{2}\left(2\pi \left( \delta(\omega_c - \omega) + \delta(\omega_c + \omega) \right) \right) \\
		&= \pi \left( \delta(\omega_c - \omega) + \delta(\omega_c + \omega) \right) \\
	\end{align*}

	\begin{step}
		Compute $Y(\omega)$.
	\end{step}

	As said earlier, since $y(t) = x(n)\cos(\omega_{c}t)$ and $X(\omega)$ and $H(\omega)$ are the frequency responses of $x(n)$ and $\cos(\omega_{c}{t})$ respectively, we know that $Y(\omega) = X(\omega) \ast H(\omega)$. Thus:

	\begin{align*}
		Y(\omega) &= \int_{-\infty}^{\infty}X(\Omega)H(\omega - \Omega)d\Omega \\
		&= \pi \int_{-\infty}^{\infty}X(\Omega)\left( \delta(\omega_{c} - (\omega - \Omega)) + \delta(\omega_{c} + (\omega - \Omega)) \right)d\Omega \\
		&= \pi \int_{-\infty}^{\infty}X(\Omega)\left( \delta(\omega_{c} - \omega + \Omega) + \delta(\omega_{c} + \omega - \Omega) \right)d\Omega \\
		&= \pi \left( X(\omega - \omega_{c}) + X(\omega + \omega_{c}) \right) \\
	\end{align*}
\end{solution}

\begin{question}[Filters]
\end{question}

\begin{subquestion}
	Compute the frequency response of a filter given by the difference equation:

	\begin{align*}
		\ddot{y} + 2\omega_{0}\dot{y} + {\omega_0}^2y &= x
	\end{align*}
\end{subquestion}

\begin{solution}
	We know that the frequency response, $H(\omega)$, of the system is given by $y(t) = H(\omega)e^{i\omega{t}}$ when the input to the filter is given by $x(t) = e^{i\omega{t}}$. Thus, differentiating, we find that:

	\begin{align*}
		\dot{y}(t) &= {i\omega}H(\omega){e^{i\omega{t}}} \\\\
		\ddot{y}(t) &= \left( i\omega \right)^2H(\omega){e^{i\omega{t}}} \\
		&= -\omega^2H(\omega) e^{i\omega{t}} \\
	\end{align*}

	Thus, by substituting, we find that:

	\begin{gather*}
		\ddot{y} + 2\omega_{0}\dot{y} + {\omega_0}^2y = x \implies \\[6pt]
		-\omega^2H(\omega)e^{i\omega{t}} + 2\omega_0{i\omega}H(\omega)e^{i\omega{t}} + {\omega_0}^2H(\omega)e^{i\omega{t}} = e^{i\omega{t}} \\[6pt]
		-\omega^2H(\omega) + 2\omega_0{i\omega}H(\omega) + {\omega_0}^2H(\omega) = 1 \\[6pt]
		H(\omega)\left( -\omega^2 + 2i\omega_0\omega + {\omega_0}^2 \right) = 1 \\[6pt]
		H(\omega) = \frac{1}{-\omega^2 + 2i\omega_0\omega + {\omega_0}^2}
	\end{gather*}
\end{solution}

\begin{subquestion}
	Design a minimum order, discrete filter that removes DC and passes the highest possible frequency with a gain of 1.
\end{subquestion}

\begin{solution}

	\begin{step}
		Construct a prototype filter.
	\end{step}

	Let $\tilde{H}(\omega)$ be the prototype filter which we will use to construct our final filter. We know that $\tilde{H}(0) = 0$ and $|H(\pi)| = 1$, since DC corresponds to a frequency of 0 and the highest possible discrete frequency is $\pi$. Thus, we find that:

	\begin{align*}
		\tilde{H}(\omega) &= \left( 1 - e^{-i\omega} \right) \\
	\end{align*}

	Substituting in $\omega = \pi$:

	\begin{align*}
		|H(\pi)| &= |1 - e^{-i\pi}| \\
		&= |1 - (-1)| \\
		&= |2| \\
		&= 2
	\end{align*}

	\begin{step}
		Construct the final filter.
	\end{step}

	Thus, normalizing our prototype filter by multiplying it by a factor of $\frac{1}{2}$, we get our final filter, $H(\omega)$, satisfying $|H(\pi)| = 1$. Thus:

	\begin{align*}
		H(\omega) &= \frac{1}{2}\left(1 - e^{-i\omega} \right) \\
	\end{align*}

	\begin{step}
		Find the difference equation of the filter.
	\end{step}

	Multiplying both sides of our final filter by $e^{i\omega{n}}$, and recognizing that $y(n) = H(\omega)e^{i\omega{n}}$ when $e^{i\omega{n}}$, we find the difference equation of the system to be:

	\begin{gather*}
		H(\omega) = \frac{1}{2}\left( 1 - e^{-i\omega} \right) \implies \\[6pt]
		H(\omega)e^{i\omega{n}} = \frac{1}{2}\left(e^{i\omega{n}} - e^{-i\omega}e^{i{\omega}{n}} \right) \\[6pt]
		y(n) = \frac{1}{2}x(n) - \frac{1}{2}x(n - 1)
	\end{gather*}
\end{solution}

\begin{question}[The Discrete Fourier Transform]
	Compute the discrete Fourier transform of the signal $x(n) = \delta(n) - \delta(n - 2)$ using $N = 4$.
\end{question}

\begin{solution}
	First, since $N = 4$, we can find that $\omega_0 = \frac{2\pi}{4} = \frac{\pi}{2}$

	The discrete Fourier transform is given by the formula:

	\begin{align*}
		X_k &= \sum_{n = 0}^{N - 1}x(n)e^{-i\omega_0{n}{k}} \\
	\end{align*}

	Thus, substituting in for $x(n)$ and $N$, we can simplify a little bit before computing the coefficients to help make things a little easier and less confusing.

	\begin{align*}
		X_k &= \sum_{n = 0}^{N - 1}x(n)e^{-i\omega_0{n}{k}} \\
		&= \sum_{n = 0}^{3}\left( \delta(n) - \delta(n - 2) \right)e^{-i\frac{\pi}{2}{n}{k}} \\
		&= \sum_{n = 0}^{3}\delta(n)e^{-i\frac{\pi}{2}{n}{k}} - \sum_{n = 0}^{3}\delta(n - 2)e^{-i\frac{\pi}{2}{n}{k}} \\
		&= 1 - e^{-2i\frac{\pi}{2}{k}} \\
	\end{align*}

	Now, we can see that:

	\begin{alignat*}{3}
		X_0 &= 1 - e^{-2i\frac{\pi}{2}(0)} &&= 1 \; \; - \; \; 1 &= 0 \\
		X_1 &= 1 - e^{-2i\frac{\pi}{2}(1)} &&= 1 - (-1) \; &= 2 \\
		X_2 &= 1 - e^{-2i\frac{\pi}{2}(2)} &&= 1  \; \; - \; \; 1 &= 0 \\
		X_3 &= 1 - e^{-2i\frac{\pi}{2}(3)} &&= 1 - (-1) \; &= 2 \\
	\end{alignat*}
\end{solution}


\end{document}
